% PDF 1.4 without object streams, for older viewers.
\pdfminorversion=4
\pdfobjcompresslevel=0
\documentclass[11pt]{article}
\usepackage[margin=2.3cm]{geometry}
\usepackage{amsmath,amssymb}
\usepackage{booktabs,array,longtable}
\usepackage[T1]{fontenc}
\usepackage{lmodern}
\usepackage{microtype}
\usepackage{tikz}
\usepackage[hidelinks]{hyperref}

\newcommand{\Sq}{\mathrm{Sq}}
\newcommand{\F}{\mathbb{F}_2}
\newcommand{\A}{\mathcal{A}}
\DeclareMathOperator{\LT}{LT}
\DeclareMathOperator{\Ext}{Ext}
\setlength{\parindent}{0pt}
\setlength{\parskip}{0.4em}

\title{A minimal resolution over the Steenrod algebra,\\ computed by Gröbner bases}
\author{Trace of the \texttt{groebner} crate (branch \texttt{groebner\_lin})}
\date{}

\begin{document}
\maketitle

\begin{abstract}
@ABSTRACT_LEAD@ the computation of the minimal free resolution
$\dots \to C_3 \to C_2 \to C_1 \to C_0 \to \F$ of the ground field over the mod $2$ Steenrod
algebra $\A$ by the \texttt{groebner} crate, a port of Weinan Lin's Gröbner basis method. The
range is homological degree $s \le @S_MAX@$ and stem $t - s \le @N_MAX@$, so internal degrees $t \le @T_MAX@$.
Every formula below is taken from the program's own data: the critical pairs it formed, the
elements it reduced, and the bases and differentials it produced. Only the bookkeeping of the
reduction steps was replayed outside the program, and the replay was checked to reproduce the
program's results exactly.
\end{abstract}

\section{Setting}

\paragraph{The algebra.}
We work in the Milnor basis $\Sq(R) = \Sq(r_1, r_2, \dots)$ of $\A$, writing $\Sq^r = \Sq(r)$.
Write $r_n = \sum_i d_{i,n} 2^i$ in binary. The digit $d_{i,n}$ records the power
$\xi_n^{2^i}$ of the dual basis, of degree $2^i(2^n - 1)$, and we record $\Sq(R)$ by the
monomial
\[ x^R = \prod_{d_{i,n} = 1} x_{in} \]
in the associated graded algebra of the May filtration
\[ E^0\A = \bigotimes_{i \ge 0,\, n \ge 1} \F[x_{in}]/(x_{in}^2). \]
For example $\Sq^3 = \Sq(3) \leftrightarrow x_{01}x_{11}$, $\Sq(0,1) \leftrightarrow x_{02}$ and
$\Sq(1,1) \leftrightarrow x_{01}x_{02}$. In $E^0\A$ a monomial is just a set of positions, so
divisibility, quotients and least common multiples are inclusion, difference and union of sets.
The positions relevant here, in increasing degree, are
@POSITIONS@.

\paragraph{The monomial order.}
A term of a free module is $\Sq(R)\,g$ with $g$ a generator. Terms are compared by
\begin{enumerate}
\item the generator: the generator added later (of higher degree) comes first;
\item the May weight $\sum_{d_{i,n} = 1} (2n - 1)$: smaller comes first;
\item the digits read as a binary number, with the position of lowest degree most significant:
  smaller comes first.
\end{enumerate}
The \emph{leading term} $\LT(x)$ of an element is its first term in this order. Elements are
always printed with their leading term first. The point of this order is that the leading term
of a product $\Sq(R)\cdot x$ is $x^R \cdot \LT(x)$ computed in $E^0\A$, \emph{whenever that
product is non-zero}. Products of leading terms vanish exactly when a position occurs twice, and
those are the cases the algorithm treats separately (``singles'' below).

\paragraph{Names.}
The generator of $C_0$ is $\iota$. Generators of $C_s$, $s \ge 1$, are named after the classes of
$\Ext_\A^{s,t}(\F, \F)$ they represent: $[h_0], [h_1], [h_2], [h_3] \in C_1$, $[h_0^2], [h_1^2],
[h_0h_2], \dots \in C_2$, $[c_0], [Ph_1]$, and so on. The names of the decomposable classes are
found from the differentials: in a minimal resolution, the product $h_i x$ of the class $x$ dual to
a generator $g$ of $C_s$ takes the value $1$ on a generator $g'$ of $C_{s+1}$ exactly when
$\Sq^{2^i} g$ occurs in $d g'$. A generator is named $h_i x$ when that product is dual to it alone.
@NAME_NOTE@ These names are only labels; the program numbers the generators of
each $C_s$ in the order they are found.

\section{The algorithm}

\paragraph{Levels.}
Level $s$ is a Gröbner basis $b^s_0, b^s_1, \dots$ of the graph of the differential,
\[ \Gamma_s = \{ (d y, y) : y \in C_{s+1} \} \subseteq C_s \oplus C_{s+1}, \]
for the order on the first component. Each $b^s_k = (x_1 \mid x_2)$ satisfies $d x_2 = x_1$,
and its leading term is that of $x_1$. The first components form a Gröbner basis of
$\operatorname{im}(d \colon C_{s+1} \to C_s)$, and pairs whose first component reduces to zero
produce elements of the kernel. The module being resolved enters as ``level $-1$'': the relations
$\Sq^{2^i}\iota = 0$ presenting $\F = \A/\A\{\Sq^1, \Sq^2, \Sq^4, \dots\}$.

\paragraph{Critical pairs and singles.}
When an element with leading term $L = x^{R}\,g$ joins level $s$, the program forms
\begin{itemize}
\item a \emph{critical pair} with each earlier element $b$ whose leading term $L'$ lies on the
  same generator: with $M = \operatorname{lcm}(L, L')$ in $E^0\A$, the S-element is
  $\Sq(M/L')\cdot b + \Sq(M/L)\cdot b_{\text{new}}$, in which the leading terms cancel. The pairs
  are pruned by the Gebauer--Möller criteria (divisibility and the chain criterion);
\item a \emph{critical single} $\Sq(x_{in}) \cdot b_{\text{new}}$ for each position $x_{in}$
  dividing $L$: then $x_{in} \cdot L = 0$ in $E^0\A$, so the leading term of the product is not
  predicted by the leading terms and the product must be reduced.
\end{itemize}
A pair or single has the degree of its leading monomial, and it is filed under that degree.

\paragraph{One degree.}
The program runs through $t = 1, 2, \dots$. In degree $t$:
\begin{enumerate}
\item Every level hands out its pairs and singles of degree $t$, and the relations of degree $t$
  are taken up.
\item Each of them is reduced: the first component modulo level $s$, carrying the second component
  along so that $d x_2 = x_1$ is preserved; then the second component modulo the first components
  of level $s+1$, i.e.\ modulo the boundaries $d(C_{s+2})$ known so far, which preserves
  $d x_2 = x_1$ because $dd = 0$. If the first component has become zero, the second component
  $z \in C_{s+1}$ is a cycle.
\item Going down from the top level, the results are compared with the elements found earlier in
  the same degree (``triangulation''), and then sorted:
  \begin{itemize}
  \item first component non-zero: a \textbf{new basis element} of level $s$;
  \item both components zero: \textbf{useless};
  \item first component zero, second component $z \ne 0$: $z$ is a cycle that is not a boundary
    (it is reduced modulo the whole image in degree $t$), so it becomes $d$ of a
    \textbf{new generator} $g$ of $C_{s+2}$, and $(z \mid g)$ joins level $s+1$. If $s + 2$
    exceeds the range, the outcome is \textbf{beyond} instead.
  \end{itemize}
\end{enumerate}
Because the levels are handled from the top down, the new elements of level $s+1$ in degree $t$
are known when the cycles of level $s$ are tested, which is what makes the resolution minimal.
No linear algebra is ever done: everything is rewriting of sparse sums of terms.

\paragraph{Truncation.}
A cycle found at level $s$ in degree $t$ is a generator of $C_{s+2}$ in stem $t-s-2$. Since we
want stems $\le @N_MAX@$, level $s$ is only needed for $t \le s + @N_MAX2@$, and pairs of higher degree are
never formed. Likewise the relations are only needed for $t \le @N_MAX1@$.

\paragraph{Running it.}\sloppy
The trace was produced by a small driver around the public API of the crate
(\texttt{Resolution::sphere} with \texttt{Bounds \{ s\_max: @S_MAX@, n\_max: @N_MAX@ \}} at $p = 2$,
\texttt{SteenrodKind::Classical}), which calls \texttt{compute\_through(t)} for $t = 1, \dots, @T_MAX@$
and records every reduction through a \texttt{PairObserver}. The driver then redid each reduction
with the crate's own products, recording the individual steps, and checked after every degree that
the outcomes, the new basis elements and the differentials agree with the program's. It also
replays the divisibility criterion on the final leading terms, and checks that it prunes as many
pairs in each degree as the program counted. The driver and the script that renders this document
are in \texttt{docs/trace}.

@STEPS@

\section{The result}

@RESULT@
